\section{Framework and Methods}

\subsection{Framework Overview}

The framework separates forecast generation from deployable planning strategy selection. Candidate forecasting models first produce forecasts for each planning unit and period. A decision layer then selects, combines, or smooths the candidate signals before converting them into executable planning signals. The resulting plans are evaluated using both forecasting metrics and operational metrics that reflect inventory exposure, plan volatility, model-switching burden, and execution-capacity violations.

This structure is designed to evaluate deployable strategy families rather than only individual forecasting models. A method can be attractive because it is accurate, because it produces stable plans, because it reduces execution burden, or because it offers a favorable balance among these objectives. The empirical design therefore reports raw metrics and normalized planning loss rather than reducing the analysis to a single forecasting leaderboard.

\subsection{Accuracy-First Methods}

Accuracy-first methods select forecasts using validation or test-period forecast accuracy. The Global Best or BestAccuracy strategy applies the numerically strongest candidate model globally. The Family Best strategy allows model choice to vary by product family or comparable planning group. These strategies are important baselines because they represent the common deployment logic of choosing the best forecaster before considering downstream execution consequences.

\subsection{Stability-First Methods}

Stability-first methods prioritize low plan volatility. In the current experiments, moving-average and BestStability variants serve as deliberately conservative references. They are not included to claim that simple smoothing is always better than machine learning. They are included because an over-stable but weak baseline helps identify the operational cost of suppressing plan changes too aggressively.

\subsection{Ensemble-Based Methods}

The Simple Ensemble averages candidate forecasts without using operational-loss information. It provides a transparent reference for forecast combination. The Operational-Loss Ensemble assigns nonnegative weights that sum to one using validation normalized planning loss rather than forecast error alone. This family tests whether soft forecast composition can preserve useful accuracy while reducing abrupt execution burden.

\subsection{Smoothing-Based Methods}

FeasibilityAwareSmoothed strategies modify the executed planning signal after a candidate plan is formed. With fixed-alpha smoothing, the adaptation rate is configured directly. With utility-CV smoothing, the adaptation rate is selected from a small validation grid using blocked validation folds to reduce dependence on a single validation period. The smoothing equation is
\begin{equation}
    x_{i,t}^{\mathrm{exec}}
    =
    \alpha_{t} x_{i,t}^{\mathrm{cand}}
    +
    (1-\alpha_{t})x_{i,t-1}^{\mathrm{exec}}.
\end{equation}
Smoothing changes the executed planning signal, not the underlying demand realization and not the raw forecast itself. It is intended to model gradual operational adaptation when execution infrastructure cannot absorb abrupt target changes.

\subsection{Feasibility-Aware Adaptive Selection}

The Feasibility-Aware Selector evaluates candidate model changes using expected planning utility rather than forecast error alone. In the default operational objective, it considers expected inventory cost, planning volatility, model-switching burden, and expected execution violations; forecast error remains a reported diagnostic unless an explicit sensitivity setting assigns it positive weight. A switch is accepted only when the expected planning utility improves after feasibility penalties. This method is intentionally interpretable; it is not a reinforcement-learning or black-box control policy.

\subsection{Finite-Horizon DP Selection}

The GreedyFeasibilitySelector minimizes the current-period expected operational cost using validation-derived model costs and current forecast-implied transition costs. It is a one-step policy and can switch when the immediate score improves.

The DPFeasibilitySelector solves the deterministic finite-horizon model-selection problem for each planning series. Its state is the previously selected model, which determines the previous planning signal and therefore the next transition cost. The DP minimizes cumulative expected operational loss over the available horizon using validation-derived expected costs; it does not use realized test demand.

The finite-horizon DP is used to quantify the value of planning ahead under switching and execution costs. Greedy selection is a one-step policy, while DP minimizes cumulative loss over the horizon. The greedy-to-DP gap therefore measures myopic selection suboptimality under the same validation-derived cost information.

\begin{center}
\fbox{\begin{minipage}{0.92\linewidth}
\textbf{Proposition 1 (Greedy selection as one-step DP).}
Let $c_t(m \mid m_{t-1}, x_{t-1})$ denote the validation-derived stage cost of choosing model $m$ at period $t$, conditional on the previous model $m_{t-1}$ and previous executed planning signal $x_{t-1}$. The finite-horizon DP value function is
\begin{equation}
V_t(m_{t-1}, x_{t-1})
=
\min_{m \in \mathcal{M}_t}
\left\{
c_t(m \mid m_{t-1}, x_{t-1})
+ V_{t+1}(m, x_t^{m})
\right\},
\quad
V_{T+1}=0.
\end{equation}
The GreedyFeasibilitySelector is equivalent to this DP when the remaining horizon is one period, or when the continuation value $V_{t+1}$ is set to zero for all states. Therefore, any performance gap between GreedyFeasibilitySelector and DPFeasibilitySelector measures myopic selection suboptimality under the same candidate forecasts and the same validation-derived stage-cost information.
\end{minipage}}
\end{center}

The BudgetedDPFeasibilitySelector adds a hard maximum switch budget $K$ to the finite-horizon DP. Unlike the soft switching penalty, this budget constrains the feasible path set. This variant is useful for planning environments where governance, monitoring, or approval processes impose a strict limit on model changes over a planning window. If missing candidate availability leaves no budget-feasible path, the implementation uses a pipeline-safe incumbent-stays fallback and marks the affected rows with fallback metadata. The fallback chooses the first period's lowest-cost model and holds that incumbent when it remains available. If the incumbent is missing later, the output is marked as potentially non-strict. These fallback cases are diagnostic and should not be interpreted as silently valid strict budgeted-DP solutions.

\subsection{Oracle Diagnostic}

The Realized-Inventory Oracle DP is non-deployable. It uses period-specific realized test inventory outcomes keyed by series, candidate model, and date. It is not a perfect-forecast oracle and does not replace model forecasts with realized demand. Instead, it is an upper-bound diagnostic for the value of knowing inventory-outcome information before deployment.

The Full-Outcome Oracle DP is an additional non-deployable benchmark over the same forecast candidates. It uses realized test-period forecast error and realized inventory outcome cost in the DP stage cost. Under the default objective, direct forecast error has zero weight, so this benchmark matters primarily for explicit forecast-error sensitivity settings. It is included to make the oracle ladder explicit without changing the deployable selector semantics.

The Realized Demand Oracle is a separate non-deployable diagnostic that uses realized demand as the forecast. Because realized demand may be more volatile than a feasible planning signal, a perfect-demand reference can have zero forecast error while still producing large execution penalties. The paper therefore treats Oracle diagnostics as information-access references rather than deployable planning strategies.

\subsection{Analysis Workflow}

Figure~\ref{fig:feasibility-analysis-flowchart} summarizes the current analysis workflow. The workflow begins with candidate forecasts, converts them to planning signals through alternative strategy families, evaluates raw operational metrics, normalizes planning loss relative to an accuracy-first reference, and then studies sensitivity to execution-risk scenarios.

\begin{figure}[htbp]
\centering
\begin{tikzpicture}[
    node distance=0.95cm and 1.05cm,
    stage/.style={
        draw,
        rounded corners=2pt,
        align=center,
        minimum width=2.7cm,
        minimum height=0.82cm,
        fill=blue!5,
        line width=0.45pt
    },
    decision/.style={
        draw,
        diamond,
        aspect=2.0,
        align=center,
        inner sep=1.2pt,
        fill=orange!10,
        line width=0.45pt
    },
    output/.style={
        draw,
        rounded corners=2pt,
        align=center,
        minimum width=2.7cm,
        minimum height=0.82cm,
        fill=green!7,
        line width=0.45pt
    },
    arrow/.style={-{Stealth[length=2.0mm]}, line width=0.45pt}
]
\node[stage] (data) {Operational\\demand data};
\node[stage, right=of data] (forecasts) {Candidate\\forecasts};
\node[stage, right=of forecasts] (signals) {Planning signal\\conversion};
\node[decision, below=of forecasts] (decision) {Feasibility-aware\\decision layer};
\node[stage, left=of decision] (capacity) {Execution capacity\\scenarios};
\node[output, right=of decision] (outputs) {Tradeoff tables\\and figures};

\draw[arrow] (data) -- (forecasts);
\draw[arrow] (forecasts) -- (signals);
\draw[arrow] (signals) -- (decision);
\draw[arrow] (capacity) -- (decision);
\draw[arrow] (decision) -- (outputs);
\end{tikzpicture}
\caption{Feasibility-aware demand-planning analysis workflow.}
\label{fig:feasibility-analysis-flowchart}
\end{figure}

